% REPLACES: main.tex lines 899-920 ("Private information makes entry state-dependent ..." to "... Online Appendix C.3.")
% No \NUM placeholders. The example numbers stand: at the declared signal vector the old
% low cost 1 already lay below gross profit at the lowest joint posterior
% (B_{r_1}(\phi_-(\mu_-)) = 1.773 by the formulas in Appendix A.5), so preparation
% decisions, orders, and entry do not change when that cost is zero. The pipeline
% must still confirm Table 3 at the new declaration.
Private information makes entry state-dependent even conditional on the
price. A high-value challenger is more likely to receive favorable
private information and therefore more likely to prepare. Competitive
pricing incorporates both conditional entry rates. The zero-cost
realization prepares at every price, after either private signal.
Entry therefore stays at or above \(\rho\) in each state, the
coefficient on public beliefs stays strictly positive, and the price still reveals the market's posterior. The buyer
combines that posterior with its own signal.
\hyperref[paper-appendix]{Appendix A} derives the posterior formulas,
state-dependent pricing, and trading bounds.

\hyperref[result-proposition-a-7]{Proposition A.7} gives sufficient
conditions for unique zero orders and entry \(\rho\) in the weak
economy, and unique full orders by investor signal with entry above
\(\rho\) in the strong economy. These conditions allow arbitrary mixed
signal-contingent orders and every unilateral deviation to an order in \([-1,1]\). These conditions allow the buyer's signal to be more accurate than the investor's.

The example has buyer accuracy 75\% and investor accuracy 70\%.
Entry rises from 0.850000 to 0.879438. A favorable private signal alone
does not justify high-cost preparation against the weak incumbent;
combined with a favorable price, it does against the strong incumbent.
Table~\ref{tab:3} reports this example with the other robustness checks.
The full accuracy grid, including economies outside the sufficient uniqueness region, appears in Online Appendix C.3.
